\documentclass[11pt]{article}
\usepackage[margin=1in]{geometry}
\usepackage{amsmath,amssymb,amsthm}
\usepackage{hyperref}
\usepackage{enumitem}

\newtheorem{theorem}{Theorem}
\newtheorem{lemma}{Lemma}
\newtheorem{proposition}{Proposition}
\newtheorem{remark}{Remark}

\newcommand{\R}{\mathbb{R}}
\newcommand{\E}{\mathbb{E}}
\newcommand{\eps}{\varepsilon}
\newcommand{\lp}{\ell_p}
\newcommand{\dimlp}{\mathrm{dim}_{\ell_p}}
\newcommand{\sgn}{\mathrm{sgn}}
\newcommand{\rad}{\varepsilon}

\title{Sharp \(p\)-Dependence for the \(\ell_p\) Margin Dimension}
\author{fa\_banach\_001, agent\_lane\_06}
\date{2026-06-15}

\begin{document}
\maketitle

\section*{Status}

This packet gives a likely valid full resolution of the \(p\)-dependence
question following Proposition 3.5 in Ashlagi--Livni--Moran--Waknine,
\emph{Margin in Abstract Spaces}, arXiv:2603.07221.

The source paper proves the \(\ell_p\) margin taxonomy
\[
 \dim_{\ell_p}(\gamma)=
 \begin{cases}
   \Theta(\gamma^{-q}), & 1<p\le 2,\quad q=p/(p-1),\\
   \Theta_p(\gamma^{-2}), & 2<p<\infty,
 \end{cases}
\]
and then asks for bounds that are also tight in their dependence on \(p\).
The result below identifies that dependence: the low-\(p\) formula is exact,
and the high-\(p\) dependence is linear in \(p\), with the sharp
\(p\gamma^{-2}\) scale in the small-margin regime.

\section*{Source Question}

The relevant source passage is Proposition 3.5 and the paragraph following it
on page 8.  For \(p>2\), the source records
\(\dim_{\ell_p}(\gamma)=\Theta_p(\gamma^{-2})\), with the hidden constant
bounded between an absolute constant and \(p\).  It then says that, although
the dependence on the margin parameter \(\gamma\) is tight in all regimes,
obtaining bounds also tight in their dependence on \(p\) remains open.

\section*{Definitions}

We use the source paper's notation.  For a Banach space \(X\),
\(\dim_X(\gamma)\) is the largest size of a \(\gamma\)-shattered subset of the
unit ball of \(X\) by norm-one linear functionals.  Corollary 3.8 of the
source paper says that \(x_1,\ldots,x_n\in B_X\) are \(\gamma\)-shattered if
and only if
\[
  \left\|\sum_{i=1}^n a_i x_i\right\|_X
  \ge \gamma\sum_{i=1}^n |a_i|
  \qquad (a_i\in\R).
\]
Equivalently, the map \(T:\ell_1^n\to X\), \(Te_i=x_i\), is an embedding with
\(\|T\|\le 1\) and \(\|T^{-1}\|\le \gamma^{-1}\) on its range.

\section*{Idea of the Proof}

The upper bounds come from testing a shattered set against random signs.  If
\(x_1,\ldots,x_n\) are \(\gamma\)-shattered, then every sign sum
\(\sum \eps_i x_i\) has norm at least \(\gamma n\).  Averaging over signs and
using scalar Khintchine inequalities gives a direct restriction on \(n\).
For \(1<p\le2\), the relevant Khintchine upper constant is \(1\), giving the
exact upper bound \(n\le\gamma^{-q}\).  For \(p>2\), the best Khintchine
constant \(B_p\) gives \(n\le B_p^2\gamma^{-2}\), and Haagerup's estimate has
\(B_p^2\asymp p\).

The lower bounds use explicit copies of \(\ell_1^n\) in finite-dimensional
\(\ell_p\).  For \(1<p\le2\), the standard basis is optimal.  For \(p>2\), the
right construction is the Rademacher system on the sign cube.  Hitczenko's
sharp moment estimate gives the missing \(p\) factor in the small-margin
regime.  A separate sign-alignment estimate shows that for every fixed
\(\gamma<1\), the high-\(p\) dimension is still linear in \(p\).  Finally,
Clarkson's inequality explains why one cannot expect a uniform lower bound
\(c p\gamma^{-2}\) up to \(\gamma=1\): there is a genuine endpoint transition.

The final push adds an exact high-\(p\) two-parameter formula in the range
\(p\ge3\).  It uses Schechtman's many-function Hanner inequality, which says
that for \(p\ge3\) the Rademacher \(p\)-moment of an \(L_p\)-valued sign sum is
dominated by the scalar sign sum of the individual \(L_p\)-norms.  The lower
bound is the Rademacher-system construction again, now sharpened by the
observation that a symmetric convex function is minimized on the \(\ell_1\)
simplex at equal coefficients.

\section*{Main Result}

\begin{theorem}[Sharp \(p\)-dependence for \(\ell_p\) margin shattering]
Let \(0<\gamma<1\).
\begin{enumerate}[label=\textup{(\alph*)}]
\item If \(1<p\le2\) and \(q=p/(p-1)\), then
\[
  \dim_{\ell_p}(\gamma)=\big\lfloor \gamma^{-q}\big\rfloor .
\]
\item If \(2<p<\infty\), then
\[
  \dim_{\ell_p}(\gamma)\le B_p^2\gamma^{-2}\le C p\,\gamma^{-2},
\]
where \(B_p\) is the optimal upper Khintchine constant and \(C\) is universal.
\item There are universal constants \(c>0\) and \(\gamma_0\in(0,1)\) such
that for all \(2<p<\infty\) and \(0<\gamma\le\gamma_0\),
\[
  \dim_{\ell_p}(\gamma)\ge c p\,\gamma^{-2}.
\]
Consequently, uniformly in \(p>2\) and small \(\gamma\),
\[
  \dim_{\ell_p}(\gamma)=\Theta(p\gamma^{-2}).
\]
\item For every \(2<p<\infty\) and \(0<\gamma<1\),
\[
  \dim_{\ell_p}(\gamma)\ge
  \max\left\{1,\left\lfloor p\log_2(1/\gamma)\right\rfloor\right\}.
\]
In particular, for every fixed \(\gamma\in(0,1)\),
\[
  \dim_{\ell_p}(\gamma)=\Theta_\gamma(p)\qquad (p\to\infty).
\]
\item If \(2<p<\infty\) and \(\gamma>2^{-1/p}\), then
\[
  \dim_{\ell_p}(\gamma)=1.
\]
\item If \(3\le p<\infty\), then the exact value is
\[
  \dim_{\ell_p}(\gamma)
  =
  \max\left\{
    n\ge1:
    \left(\E\left|\frac{\rad_1+\cdots+\rad_n}{n}\right|^p\right)^{1/p}
    \ge \gamma
  \right\}.
\]
\end{enumerate}
Thus the \(p\)-dependence in the source paper's \(\ell_p\) taxonomy is sharp:
there is no hidden \(p\)-loss for \(1<p\le2\), and the high-\(p\) dependence is
linear in \(p\).
\end{theorem}

\section*{Proof}

\begin{proof}

We first prove the upper bounds.  Let \(x_1,\ldots,x_n\) be
\(\gamma\)-shattered in the unit ball of \(\ell_p\).  By the source
characterization, for every sign vector \(\eps\in\{-1,1\}^n\),
\[
  \left\|\sum_{i=1}^n \eps_i x_i\right\|_p\ge \gamma n.
\]

Suppose first that \(1<p\le2\).  By the scalar Khintchine inequality with
upper constant \(1\) for \(p\le2\),
\[
 \E_\eps \left|\sum_{i=1}^n \eps_i a_i\right|^p
 \le \left(\sum_{i=1}^n |a_i|^2\right)^{p/2}
 \le \sum_{i=1}^n |a_i|^p .
\]
Applying this coordinatewise gives
\[
(\gamma n)^p
\le
\E_\eps\left\|\sum_{i=1}^n \eps_i x_i\right\|_p^p
=
\sum_j \E_\eps\left|\sum_{i=1}^n \eps_i x_i(j)\right|^p
\le
\sum_{i=1}^n \|x_i\|_p^p
\le n.
\]
Thus \(n\le\gamma^{-p/(p-1)}=\gamma^{-q}\).

The reverse inequality for \(1<p\le2\) is supplied by the standard basis.
Indeed, if \(e_1,\ldots,e_n\) are the first \(n\) unit vectors and
\(\sum_i |a_i|=1\), then
\[
  \left\|\sum_{i=1}^n a_i e_i\right\|_p
  =\|a\|_p
  \ge n^{1/p-1}\|a\|_1
  = n^{-1/q}.
\]
Hence the standard basis is \(\gamma\)-shattered whenever
\(n\le\gamma^{-q}\), proving
\(\dim_{\ell_p}(\gamma)=\lfloor\gamma^{-q}\rfloor\).

Now let \(p>2\).  The same sign test, together with the scalar Khintchine
upper constant \(B_p\), gives
\[
\begin{aligned}
(\gamma n)^p
&\le
\E_\eps\left\|\sum_{i=1}^n \eps_i x_i\right\|_p^p \\
&=
\sum_j \E_\eps\left|\sum_{i=1}^n \eps_i x_i(j)\right|^p \\
&\le
B_p^p \sum_j \left(\sum_{i=1}^n |x_i(j)|^2\right)^{p/2}.
\end{aligned}
\]
Since \(p/2\ge1\),
\[
 \left(\sum_{i=1}^n |x_i(j)|^2\right)^{p/2}
 \le n^{p/2-1}\sum_{i=1}^n |x_i(j)|^p.
\]
Summing in \(j\) and using \(\|x_i\|_p\le1\), we get
\[
(\gamma n)^p\le B_p^p n^{p/2}.
\]
Therefore \(n\le B_p^2\gamma^{-2}\).  Haagerup's sharp Khintchine constants
give \(B_p^2\le C p\), proving the asserted high-\(p\) upper bound.

We next prove the small-margin high-\(p\) lower bound.  We use Hitczenko's
Rademacher moment estimate in the following standard form: there is a
universal \(A>0\) such that, for every \(p\ge2\), every finite scalar sequence
\(b=(b_i)\), and independent Rademacher signs \((\eps_i)\), if
\((b_i^*)\) is the non-increasing rearrangement of \((|b_i|)\), then
\[
 \left(\E\left|\sum_i b_i\eps_i\right|^p\right)^{1/p}
 \ge
 A\left(
   \sum_{i\le \lfloor p\rfloor} b_i^*
   + \sqrt p
       \left(\sum_{i>\lfloor p\rfloor} (b_i^*)^2\right)^{1/2}
 \right).
\]
Consequently, if \(b\in\R^n\) and \(\sum_i |b_i|=1\), then
\[
 \left(\E\left|\sum_{i=1}^n b_i\eps_i\right|^p\right)^{1/p}
 \ge
 c_0\min\left(1,\sqrt{\frac{p}{n}}\right)
\]
with a universal \(c_0>0\).  Indeed, if \(n\le\lfloor p\rfloor\), the head
term is \(1\).  If \(n>\lfloor p\rfloor\), either the head has mass at least
\(1/2\), or the tail has \(\ell_1\)-mass at least \(1/2\), in which case its
\(\ell_2\)-norm is at least \(1/(2\sqrt n)\).

Let \(n=\lfloor c_1 p\gamma^{-2}\rfloor\), with \(c_1>0\) sufficiently small,
and let \(\Omega=\{-1,1\}^n\) with the uniform probability measure.  For
\(1\le i\le n\), let \(r_i(\omega)=\omega_i\).  Embed \(L_p(\Omega)\)
isometrically into \(\ell_p^{2^n}\) by
\[
 f\mapsto 2^{-n/p}(f(\omega))_{\omega\in\Omega},
\]
and let \(x_i\) be the image of \(r_i\).  Then \(\|x_i\|_p=1\), and for
\(\sum_i |a_i|=1\),
\[
 \left\|\sum_{i=1}^n a_i x_i\right\|_p
 =
 \left(\E_\omega\left|\sum_{i=1}^n a_i r_i(\omega)\right|^p\right)^{1/p}
 \ge c_0\min\left(1,\sqrt{\frac{p}{n}}\right).
\]
For \(0<\gamma\le\gamma_0\), with \(\gamma_0\) universal and \(c_1\) chosen
small enough, the last expression is at least \(\gamma\).  Thus the
\(x_i\)'s are \(\gamma\)-shattered and
\(\dim_{\ell_p}(\gamma)\ge c p\gamma^{-2}\).

For the fixed-margin high-\(p\) lower bound, use the same Rademacher system
but choose \(n\le p\log_2(1/\gamma)\).  For every \(a\in\R^n\) with
\(\sum_i |a_i|=1\), there is a sign vector \(\omega\) such that
\(\sum_i a_i r_i(\omega)=1\), namely \(\omega_i=\sgn(a_i)\) on the support of
\(a\).  Hence
\[
 \left(\E_\omega\left|\sum_{i=1}^n a_i r_i(\omega)\right|^p\right)^{1/p}
 \ge 2^{-n/p}.
\]
If \(n\le p\log_2(1/\gamma)\), this is at least \(\gamma\), so the
Rademacher vectors are \(\gamma\)-shattered.  This proves
\[
 \dim_{\ell_p}(\gamma)\ge
 \max\left\{1,\left\lfloor p\log_2(1/\gamma)\right\rfloor\right\}.
\]
Together with the upper bound \(O(p\gamma^{-2})\), it follows that, for every
fixed \(\gamma<1\), \(\dim_{\ell_p}(\gamma)=\Theta_\gamma(p)\).

It remains to justify the endpoint obstruction.  Suppose \(p>2\) and two unit
vectors \(x,y\in\ell_p\) are \(\gamma\)-shattered.  Applying the source
characterization to \((x,y)\) with coefficients \((1,1)\) and \((1,-1)\)
gives
\[
  \left\|\frac{x+y}{2}\right\|_p\ge\gamma,
  \qquad
  \left\|\frac{x-y}{2}\right\|_p\ge\gamma.
\]
Clarkson's inequality for \(L_p\), \(p\ge2\), gives
\[
 \left\|\frac{x+y}{2}\right\|_p^p
 +
 \left\|\frac{x-y}{2}\right\|_p^p
 \le
 \frac{\|x\|_p^p+\|y\|_p^p}{2}
 \le1.
\]
Therefore \(2\gamma^p\le1\), i.e. \(\gamma\le2^{-1/p}\).  If
\(\gamma>2^{-1/p}\), no two-point shattered set exists.  Since one unit
vector is \(\gamma\)-shattered for every \(\gamma<1\), we have
\(\dim_{\ell_p}(\gamma)=1\) in this endpoint range.

We finally prove the exact formula for \(p\ge3\).  Let
\[
  \rho_p(n)
  =
  \left(\E\left|\frac{\rad_1+\cdots+\rad_n}{n}\right|^p\right)^{1/p}.
\]
For the lower bound, use the Rademacher coordinate functions
\(r_1,\ldots,r_n\) on \(\{-1,1\}^n\), embedded isometrically into a finite
\(\ell_p\) as above.  If \(\sum_i |a_i|=1\), signs may be absorbed into the
Rademacher variables, so assume \(a_i\ge0\).  The function
\[
  (a_1,\ldots,a_n)\mapsto
  \E\left|\sum_{i=1}^n a_i\rad_i\right|^p
\]
is convex and symmetric on \(\R^n\).  Hence its restriction to the probability
simplex is Schur-convex.  Since every point of the simplex majorizes
\((1/n,\ldots,1/n)\), this function is minimized at equal coefficients.
Therefore
\[
 \left\|\sum_{i=1}^n a_i r_i\right\|_{L_p}
 =
 \left(\E\left|\sum_{i=1}^n a_i\rad_i\right|^p\right)^{1/p}
 \ge \rho_p(n).
\]
Thus the Rademacher vectors are \(\gamma\)-shattered whenever
\(\rho_p(n)\ge\gamma\).

Conversely, suppose \(x_1,\ldots,x_n\in B_{\ell_p}\) are
\(\gamma\)-shattered and \(p\ge3\).  Schechtman's many-function Hanner
inequality gives, for \(L_p\)-valued functions \(f_i\),
\[
  \E_\rad\left\|\sum_{i=1}^n \rad_i f_i\right\|_p^p
  \le
  \E_\rad\left|\sum_{i=1}^n \rad_i\|f_i\|_p\right|^p .
\]
Applying this to \(x_i\), and then using \(\|x_i\|_p\le1\), gives the next
bound by the elementary monotonicity of scalar Rademacher moments in each
coefficient:
\[
  \E_\rad\left\|\sum_{i=1}^n \rad_i x_i\right\|_p^p
  \le
  \E_\rad|\rad_1+\cdots+\rad_n|^p .
\]
Shattering gives
\(\|\sum_i \rad_i x_i\|_p\ge\gamma n\) for every sign vector, so
\[
  (\gamma n)^p
  \le
  \E|\rad_1+\cdots+\rad_n|^p.
\]
Equivalently, \(\rho_p(n)\ge\gamma\).  This proves the exact formula for
\(p\ge3\).
\end{proof}

\section*{Interpretation}

This theorem resolves the \(p\)-dependence asked for after Proposition 3.5.
For \(1<p\le2\), the exact answer is \(\lfloor\gamma^{-q}\rfloor\), so the
source's \(p\)-dependence through \(q=p/(p-1)\) has no additional hidden
constant.  For \(p>2\), the upper constant is \(B_p^2\asymp p\), and the
Rademacher-Hitczenko lower bound shows that this linear growth in \(p\) is
attained in the small-margin regime where the source's
\(\Theta_p(\gamma^{-2})\) asymptotic is meaningful.  The sign-alignment lower
bound further shows that linear \(p\)-growth persists for every fixed
\(\gamma<1\).

The Clarkson endpoint statement is included to prevent an overinterpretation:
one cannot have a universal lower bound \(c p\gamma^{-2}\) for every
\(\gamma\in(0,1)\), because the dimension drops to \(1\) when
\(\gamma>2^{-1/p}\).  Thus the remaining near-\(\gamma=1\) transition is a
different two-parameter endpoint phenomenon, not a failure of the
p-dependence resolution requested in the source remark.

In fact, for \(p\ge3\) the endpoint transition is completely described by the
Rademacher random-walk moment \(\rho_p(n)\).  This is sharper than needed for
the source p-dependence question.  Jenkins--Tkocz note that the corresponding
real Rademacher many-function Hanner inequality is open in the range
\(2<p<3\).  Conditional on that open inequality, the same proof gives the
same exact formula for every \(2\le p<\infty\).

\section*{Novelty and Literature Search}

Local run indexes and parsed sources were searched for arXiv:2603.07221, the
source title, \(\ell_p\) margin phrases, p-dependence phrases, fat-shattering
phrases, Rademacher moment phrases, Hitczenko, Balcan--Berlind, and
\(\ell_1^n\) embeddings into \(L_p\).  Web searches on 2026-06-15 covered the
source title and p-dependence phrase, Hitczenko/Rademacher moment estimates,
Khintchine constants, Clarkson inequalities, and Banach--Mazur or
\(\ell_1^n\)-into-\(L_p\) formulations.  These searches found the source
paper and the standard tools cited below.  The final push also identified
Schechtman's many-function Hanner inequality for real \(L_p\), \(p\ge3\), via
Jenkins--Tkocz.  This appears not to have been applied to the exact margin
dimension question in the source paper.

\section*{References}

\begin{enumerate}[label={[\arabic*]}]
\item Y. Ashlagi, R. Livni, S. Moran, and T. Waknine,
\emph{Margin in Abstract Spaces}, arXiv:2603.07221, 2026.
\item U. Haagerup, \emph{The best constants in the Khintchine inequality},
Studia Mathematica 70 (1981), 231--283.
\item P. Hitczenko, \emph{On the Rademacher Series}, Probability in Banach
Spaces 9, 31--36, 1994.
\item P. Hitczenko and S. Montgomery-Smith,
\emph{Measuring the magnitude of sums of independent random variables},
Annals of Probability 29 (2001), 447--466; arXiv:math/9909054.
\item M.-F. Balcan and C. Berlind,
\emph{A New Perspective on Learning Linear Separators with Large \(L_qL_p\)
Margins}, AISTATS 2014.
\item J. A. Clarkson, \emph{Uniformly convex spaces}, Transactions of the
American Mathematical Society 40 (1936), 396--414.
\item G. Schechtman, \emph{Two remarks on 1-unconditional basic sequences in
\(L_p\), \(3\le p<\infty\)}, Geometric aspects of functional analysis
(Israel, 1992--1994), 251--254, Operator Theory: Advances and Applications 77,
Birkh\"auser, Basel, 1995.
\item J. Jenkins and T. Tkocz, \emph{Complex Hanner's inequality for many
functions}, arXiv:2207.09122, 2022.
\end{enumerate}

\section*{Human Review Recommendation}

Review as a likely valid full solution to the p-dependence question following
Proposition 3.5 of the source paper.  The main checks are the normalization in
Corollary 3.8, the coordinatewise Khintchine upper-bound computations, and the
scope interpretation: the theorem solves the p-growth problem, while the
\(p\ge3\) refinement gives an exact two-parameter endpoint formula.  The
remaining exact-formula range \(2<p<3\) is tied to the open real Rademacher
many-function Hanner inequality identified by Jenkins--Tkocz.

\end{document}
